\chapter{Calcul distribué et ordonnanceurs}\label{ch:pl:distributed-compute-and-schedulers}

Un \omterm{def:pl:distributed-compute-and-schedulers:sweep}{balayage de paramètres} de dix mille backtests a été soumis à 18~h un vendredi et s'est terminé à 13~h~47 le samedi. À partir de 10~h~20 ce matin-là,
moins d'un dixième des cœurs de la grappe étaient occupés. Quarante des tâches duraient trente fois plus que la médiane, et elles avaient
été soumises en dernier~: la grappe les a atteintes après douze heures, puis a attendu encore sept heures et demie. Rien n'avait échoué~; l'ordonnanceur avait fait
exactement ce qu'on lui avait demandé. Ce chapitre porte sur ce qu'il faut lui demander~: comment un balayage devient un \omterm{def:pl:distributed-compute-and-schedulers:graph}{graphe de tâches}, comment les ordonnanceurs ordonnent ce graphe sur
une grappe, pourquoi les dernières tâches décident de la fin d'un balayage, et ce que coûtent et permettent la mise en cache, les nouvelles tentatives et les copies de secours.

\section{Grappes de calcul et leur usage}

\begin{definition}[Grappe de calcul, ordonnanceur de tâches]\label{def:pl:distributed-compute-and-schedulers:cluster}
Une \emph{grappe de calcul}\index{grappe de calcul} est un ensemble de machines (nœuds), chacune dotée de cœurs, de mémoire et parfois d'accélérateurs, gérées
comme un seul réservoir de capacité. Un \emph{ordonnanceur de tâches}\index{ordonnanceur de tâches} décide quel travail soumis s'exécute sur quelle partie du réservoir et à quel moment,
selon une politique, les ressources que demande chaque travail et les priorités des utilisateurs qui l'ont soumis.
\end{definition}

Une plateforme de recherche utilise sa grappe pour des travaux volumineux, répétitifs et surtout indépendants~: \omterm{def:pl:distributed-compute-and-schedulers:sweep}{balayages de paramètres}, recalibrages walk-forward
(livre~7, chapitre~20), simulations sur des graines (chapitres~11 et~12), scénarios de risque (chapitres~18 et~20), entraînement de modèles (livre~12). La plupart de ces travaux
relèvent de la forme de parallélisme la plus simple qui soit.

\begin{definition}[Charge trivialement parallèle, balayage de paramètres]\label{def:pl:distributed-compute-and-schedulers:sweep}
Une charge de travail est \emph{trivialement parallèle}\index{charge trivialement parallèle} lorsqu'elle se découpe en tâches qui ne communiquent pas
pendant leur exécution. Un \emph{balayage de paramètres}\index{balayage de paramètres} est une telle charge~: le même calcul, typiquement un backtest, exécuté une fois
pour chaque point d'une grille de paramètres, de tranches de données ou de graines.
\end{definition}

Être \omterm{def:pl:distributed-compute-and-schedulers:sweep}{trivialement parallèle} ne signifie pas être trivial à terminer à temps. Les tâches d'un balayage ont des durées inégales --- une fenêtre d'historique plus longue, un marché plus actif,
une tranche de données plus lente --- et le balayage ne se termine qu'avec sa dernière tâche.

\begin{definition}[Durée totale, traînard]\label{def:pl:distributed-compute-and-schedulers:makespan}
La \emph{durée totale}\index{durée totale} d'un ensemble de tâches est le temps écoulé entre le début de la première et la fin de la dernière. Un
\emph{traînard}\index{traînard} est une tâche qui se termine bien après la plupart de ses pairs, parce qu'elle est intrinsèquement longue, parce qu'elle tourne
sur une machine lente ou surchargée, ou parce qu'elle a échoué puis a été relancée.
\end{definition}

\section{Ordonnanceurs et files d'attente}

L'ordonnanceur le plus simple tient une seule file d'attente et confie le cœur libre suivant à la tâche en tête~: premier entré, premier sorti. Graham (1969) a analysé cet
\emph{list scheduling} et en a borné la performance~: quel que soit l'ordre, la \omterm{def:pl:distributed-compute-and-schedulers:makespan}{durée totale} sur $m$ cœurs identiques est au plus $2 - 1/m$ fois la durée optimale,
et si l'on prend les tâches les plus longues en premier --- \emph{longest processing time}, LPT --- la borne se resserre à $4/3 - 1/(3m)$. Ces bornes reposent sur une
observation simple, qui sert aussi d'étalon de mesure à ce chapitre.

\begin{proposition}[Borne inférieure de la durée totale]\label{prop:pl:distributed-compute-and-schedulers:bound}
Aucun ordonnancement de tâches de durées $d_i$ sur $m$ cœurs ne se termine avant $\max\bigl(\sum_i d_i / m,\ \max_i d_i,\ L\bigr)$, où $L$ est la
plus longue chaîne de tâches dépendantes~: le travail doit être effectué par $m$ cœurs, la tâche la plus longue tourne sur un seul cœur, et une chaîne s'exécute en séquence.
\end{proposition}

\omcode{code/firm/jobgraph/firm_jobgraph.py}{279}{292}{La borne inférieure de la \cref{prop:pl:distributed-compute-and-schedulers:bound}, avec la
plus longue chaîne calculée sur le graphe de tâches.}\label{lst:pl:distributed-compute-and-schedulers:bound}

LPT suppose de savoir quelles tâches sont longues. Un ordonnanceur ne connaît pas les durées, seulement des estimations --- à partir des paramètres (une fenêtre plus longue, un historique plus long), de
précédentes exécutions du même code, ou de l'utilisateur. Les estimations du chapitre sont justes à un facteur $e^{0.3z}$ près, avec
$z$ de loi normale centrée réduite, une dispersion typique pour un modèle de temps d'exécution.

\begin{definition}[Ordonnancement à parts équitables]\label{def:pl:distributed-compute-and-schedulers:fairshare}
L'\emph{ordonnancement à parts équitables}\index{ordonnancement à parts équitables} (fair share) attribue à chaque utilisateur ou équipe une part cible de la grappe et ordonne la file pour que
ceux qui consomment moins que leur part soient servis en premier~; une équipe qui soumet dix mille tâches n'empêche pas une équipe qui en soumet cinq cents
une heure plus tard d'obtenir des cœurs.
\end{definition}

\omcode{code/firm/jobgraph/firm_jobgraph.py}{109}{129}{Partage équitable sous sa forme la plus simple~: la tâche suivante vient de l'équipe qui utilise le moins
de cœurs par unité de sa cible.}\label{lst:pl:distributed-compute-and-schedulers:fairshare}

\begin{dated}{2026-09}{Partage équitable dans un ordonnanceur de production}\label{dat:pl:distributed-compute-and-schedulers:slurm}
Slurm, un gestionnaire de charge de travail libre et largement utilisé, calcule un facteur de partage équitable à partir de parts normalisées $S$ et d'usages normalisés $U$
(«~Fair Tree~», par défaut depuis Slurm 19.05)~: le partage équitable de niveau d'une association vaut $S/U$, supérieur à 1 lorsqu'elle est sous-servie,
et les utilisateurs sont classés de sorte que les enfants d'un compte mieux servi se classent toujours derrière ceux d'un compte moins bien servi.
\end{dated}

\begin{definition}[Vol de tâches]\label{def:pl:distributed-compute-and-schedulers:stealing}
Le \emph{vol de tâches}\index{vol de tâches} (work stealing) est une politique d'ordonnancement décentralisée~: chaque travailleur tient sa propre file de tâches et prend dans sa
tête~; un travailleur dont la file est vide prend une tâche à l'autre extrémité de la file d'un autre travailleur. Blumofe et Leiserson (1999) l'ont analysé
pour des calculs multithreadés~; il équilibre la charge sans file centrale.
\end{definition}

\omcode{code/firm/jobgraph/firm_jobgraph.py}{132}{154}{Vol de tâches~: tâches réparties à tour de rôle entre des files par cœur, un cœur inactif volant
dans la plus pleine.}\label{lst:pl:distributed-compute-and-schedulers:stealing}

\section{Graphes de tâches}

\begin{definition}[Graphe de tâches]\label{def:pl:distributed-compute-and-schedulers:graph}
Un \emph{graphe de tâches}\index{graphe de tâches} est un graphe orienté sans cycle dont les nœuds sont des tâches, chacune avec une demande de ressources et une
durée estimée, et dont les arêtes indiquent quelle sortie une tâche lit d'une autre~; une tâche est prête lorsque tous ses prédécesseurs se sont terminés.
\end{definition}

Un balayage n'est presque jamais une liste plate. Ses backtests lisent des variables~; les variables lisent les données nettoyées~; les données ont été construites une fois pour de nombreux balayages.
Écrite sous forme de graphe, chaque étape partagée devient une tâche dont dépendent de nombreux backtests, et sa sortie peut être mise en cache sous la clé de contenu du
flux de recherche du livre~7 (chapitre~29)~: le hachage de son code, de ses paramètres et de ses entrées. Le balayage suivant sur les mêmes données retrouve les variables
dans le cache et n'exécute que les backtests.

\section{Balayages de paramètres et traînards}

\begin{omfigure}
\begin{tikzpicture}[font=\scriptsize,>=Stealth]
\foreach \i/\l in {0/{étape de variables 1},1/{étape de variables 2},2/{\dots},3/{étape de variables 100}} {
  \node[draw,rounded corners=2pt,fill=omMeth!15,minimum width=2.2cm] (f\i) at (\i*3,1.4) {\l};}
\foreach \i in {0,1,3} {\foreach \j in {-1,0,1} {\draw[->] (f\i.south) -- ++(\j*0.8,-0.8);}}
\foreach \i in {0,1,3} {\node[font=\tiny] at (\i*3,0.1) {100 backtests};}
\node[draw,rounded corners=2pt,fill=gray!10,align=center,minimum width=10cm] (c) at (4.5,-0.6) {grappe~: 24 nœuds $\times$ 16 cœurs~;
un nœud quatre fois plus lent~; les travaux en cours échouent à 0.002 par cœur-heure};
\end{tikzpicture}

\omcaption{Le \omterm{def:pl:distributed-compute-and-schedulers:graph}{graphe de tâches} et la grappe du chapitre. Cent étapes de variables de vingt minutes chacune, lues chacune par cent backtests~; la durée de chaque backtest est tirée
de la dispersion mesurée, quarante d'entre eux valant trente fois la médiane dans le balayage à plat. Les étapes de variables sont mises en cache
par clé de contenu.}\label{fig:pl:distributed-compute-and-schedulers:graph}
\end{omfigure}

La grappe du chapitre est simulée (\texttt{firm.jobgraph})~; ses tâches ne sont pas inventées. Vingt-quatre petits backtests du chapitre~11 ---
rejeux de niveau~3 de sessions de cinq minutes de \texttt{firm.tape} --- ont été exécutés sur cet ordinateur portable par la même interface d'ordonnancement, avec un seul
travailleur~: leur médiane valait 0.089~seconde, la dispersion de leurs logarithmes de durées 0.76, et la plus longue a pris 8.0 fois la médiane. Le balayage simulé
conserve cette dispersion et ramène la médiane à quinze minutes~: 10\,000 backtests, quarante d'entre eux valant trente fois la médiane (sept heures et demie,
une longue fenêtre sur un long historique) soumis en dernier, comme un balayage sur une grille soumet naturellement ses plus grands paramètres. Le travail
représente 3\,649 cœur-heures sur 384 cœurs, soit une borne inférieure de 9.50~heures.

\begin{omfigure}
\begin{tikzpicture}
\begin{axis}[width=12cm,height=5.6cm,xlabel={heures après la soumission},ylabel={cœurs occupés (échelle log)},xmin=0,xmax=21,
  ymode=log,ymin=0.7,ymax=600,unbounded coords=jump,
  grid=major,grid style={gray!20},tick label style={font=\small},label style={font=\small},table/col sep=comma]
\addplot[color=omThm,very thick,const plot] table[x=t_h,y=busy]{figdata/platforms/13-distributed-compute-and-schedulers/profile_fifo.csv};
\addplot[color=omDef,thick,densely dashed,const plot] table[x=t_h,y=busy]{figdata/platforms/13-distributed-compute-and-schedulers/profile_lpt.csv};
\node[omThm,font=\scriptsize,anchor=south] at (axis cs:14,45) {premier entré, premier sorti};
\node[omDef,font=\scriptsize,anchor=north east,align=left] at (axis cs:19.8,500) {plus long d'abord,\\avec copies de secours~:\\se termine à 10.0};
\draw[black,dotted] (axis cs:9.50,0.7) -- (axis cs:9.50,600) node[pos=0.3,right,font=\scriptsize]{borne inférieure};
\end{axis}
\end{tikzpicture}

\omcaption{Cœurs occupés au fil du temps pour le même balayage selon deux ordonnancements, relevés toutes les dix minutes, en échelle logarithmique. Le premier entré, premier sorti garde la grappe pleine
pendant environ neuf heures, n'atteint les quarante tâches longues qu'après douze heures, et les exécute sur une grappe presque vide jusqu'à ce que la dernière se termine, à 19.8~heures~; le plus long d'abord, avec copies de secours, se termine à 10.0~heures.
Données~: \texttt{fig\_jobgraph.py}.}\label{fig:pl:distributed-compute-and-schedulers:profile}
\end{omfigure}

La \Cref{tab:pl:distributed-compute-and-schedulers:policies} donne le balayage sous chaque politique. Le premier entré, premier sorti n'atteint la dernière des quarante
tâches longues qu'après douze heures, et le balayage se termine à 19.8~heures~--- 475.0 heures-nœud d'une grappe réservée pour un travail qui demandait 3\,649
cœur-heures. Le plus long d'abord lance les tâches longues dès le début et se termine à 11.8~heures.
Le \omterm{def:pl:distributed-compute-and-schedulers:stealing}{vol de tâches} (work stealing), qui équilibre la charge sans la réordonner, se termine à 16.3.

\begin{table}[ht]
\centering\small
\begin{tabular}{@{}>{\raggedright\arraybackslash}p{2.2cm}>{\centering\arraybackslash}p{2.3cm}>{\centering\arraybackslash}p{2.3cm}>{\centering\arraybackslash}p{2.3cm}>{\centering\arraybackslash}p{2.3cm}>{\centering\arraybackslash}p{2.3cm}@{}}
\toprule
politique & \multicolumn{2}{c}{\omterm{def:pl:distributed-compute-and-schedulers:makespan}{durée totale} (makespan, heures)} & \multicolumn{2}{c}{heures-nœud} & copies de secours \\
 & simple & avec copies de secours & simple & avec copies de secours & \\
\midrule
premier entré, premier sorti & 19.79 & 20.32 & 475.0 & 487.6 & 152 \\
le plus long d'abord & 11.80 & 10.02 & 283.3 & 240.4 & 105 \\
\omterm{def:pl:distributed-compute-and-schedulers:stealing}{vol de tâches} & 16.33 & 16.40 & 391.9 & 393.5 & 145 \\
le plus long d'abord, sans nœud lent & 10.32 & --- & 247.6 & --- & --- \\
\midrule
borne inférieure & \multicolumn{2}{c}{9.50} & & & \\
\bottomrule
\end{tabular}
\caption{Le balayage de 10\,000 tâches sur 384 cœurs sous chaque politique, sans et avec copies de secours spéculatives (une copie de secours lancée après 2.5 fois l'estimation),
et le plus long d'abord sur une grappe sans le nœud lent. Les heures-nœud sont les 24 nœuds de la grappe réservés pendant la \omterm{def:pl:distributed-compute-and-schedulers:makespan}{durée totale}.}
\label{tab:pl:distributed-compute-and-schedulers:policies}
\end{table}

\begin{omfigure}
\begin{tikzpicture}
\begin{axis}[width=12cm,height=5.4cm,ybar,bar width=12pt,ylabel={durée totale (heures)},ymin=0,ymax=26,
  symbolic x coords={FIFO,LPT,work stealing},xtick=data,xticklabels={FIFO,LPT,vol de tâches},enlarge x limits=0.25,
  tick label style={font=\small},label style={font=\small},grid=major,grid style={gray!20},table/col sep=comma,
  legend style={font=\footnotesize,at={(0.5,1.03)},anchor=south,legend columns=2},nodes near coords,
  nodes near coords style={font=\scriptsize,rotate=90,anchor=west,/pgf/number format/fixed,/pgf/number format/fixed zerofill,/pgf/number format/precision=2},point meta=y,ybar=4pt]
\addplot[fill=omThm!45,draw=omThm] table[x=policy,y=plain]{figdata/platforms/13-distributed-compute-and-schedulers/makespans.csv};
\addlegendentry{simple}
\addplot[fill=omDef!45,draw=omDef] table[x=policy,y=backups]{figdata/platforms/13-distributed-compute-and-schedulers/makespans.csv};
\addlegendentry{avec copies de secours}
\draw[black,dashed] ({rel axis cs:0,0}|-{axis cs:FIFO,9.50}) -- ({rel axis cs:1,0}|-{axis cs:FIFO,9.50})
  node[pos=0.335,below,font=\tiny]{borne inférieure 9.50};
\end{axis}
\end{tikzpicture}

\omcaption{\omterm{def:pl:distributed-compute-and-schedulers:makespan}{Durée totale} du balayage par politique, sans et avec copies de secours spéculatives (après 2.5 fois l'estimation), comparée à la borne inférieure.
Le premier entré, premier sorti (FIFO) et le \omterm{def:pl:distributed-compute-and-schedulers:stealing}{vol de tâches} exécutent les tâches longues en dernier~; le plus long d'abord (LPT) les exécute en premier. Données~:
\texttt{fig\_jobgraph.py}.}\label{fig:pl:distributed-compute-and-schedulers:makespans}
\end{omfigure}

\begin{definition}[Exécution spéculative]\label{def:pl:distributed-compute-and-schedulers:spec}
L'\emph{exécution spéculative}\index{exécution spéculative} lance une seconde copie d'une tâche qui tourne bien plus longtemps que prévu, sur
une autre machine~; la première copie à se terminer est conservée, et l'autre est interrompue.
\end{definition}

L'article MapReduce de Google a nommé le problème et son remède~: un \omterm{def:pl:distributed-compute-and-schedulers:makespan}{traînard} (straggler) peut être une machine dont le disque défaillant voit son débit de lecture tomber de 30 à
1~Mo/s, et lorsqu'un travail approche de son terme, le maître planifie des exécutions de secours pour les tâches restantes encore en cours~; leur tri a duré
44~\% de plus sans copies de secours. Le seuil compte. Les copies de secours du chapitre démarrent quand une tâche a tourné 2.5 fois son estimation~: avec
des estimations justes à $e^{0.3z}$ près, une tâche normale dépasse ce point avec une probabilité de 0.11~\% ($z > 3.05$), si bien que les copies de secours visent surtout les tâches
sur le nœud lent. Le plus long d'abord gagne 1.8~heures grâce à elles~; le premier entré, premier sorti en perd 0.5, car ses tâches longues tardent à démarrer, et non
à s'exécuter lentement, et les copies de secours ne font que prendre des cœurs~; le travail gaspillé --- copies interrompues, tentatives échouées puis relancées --- est de 91 à 130
cœur-heures, soit 2.5 à 3.6~\% du balayage.

\omcode{code/firm/jobgraph/firm_jobgraph.py}{206}{227}{Exécution spéculative dans le simulateur~: la copie de secours d'une tâche en retard est placée sur un cœur libre situé
sur un autre nœud, avant que la file ne soit servie.}\label{lst:pl:distributed-compute-and-schedulers:spec}

\begin{example}[Le balayage terminé le samedi]\label{ex:pl:distributed-compute-and-schedulers:sweep}
Le plus long d'abord avec copies de secours se termine à 10.02~heures, 5.5~\% au-dessus de la borne inférieure de 9.50~heures. Sans le nœud lent, mais aussi sans copies de secours,
il se termine à 10.32~: les copies de secours rattrapent plus que ce que coûte le nœud lent, car elles sauvent aussi les tâches dont les estimations étaient
fausses et celles qui ont échoué. Le premier entré, premier sorti sur la même grappe se termine à 19.79~heures, soit deux fois plus longtemps, et occupe 475 heures-nœud contre
240 heures-nœud. La grappe est restée presque vide pendant trois heures et demie~: à partir de 16.3~heures après la soumission, moins d'un dixième de ses cœurs étaient
occupés.
\end{example}

\section{Parts équitables entre équipes}

Une deuxième équipe soumet 500 backtests une heure après le balayage de la première. Sous premier entré, premier sorti, ils attendent derrière les 10\,000 tâches~: la moitié
d'entre eux est terminée 8.28~heures après la soumission, et la dernière 11.4~heures après. Avec un partage équitable à cibles égales, chaque cœur libéré revient à la deuxième équipe
jusqu'à ce qu'elle occupe la moitié de la grappe~: la moitié de ses tâches sont terminées en 0.59~heures et toutes en 2.64, tandis que le balayage de la première équipe se termine à
16.68~heures au lieu de 19.79 --- plus tôt, et non plus tard, car l'entrelacement des tâches de la deuxième équipe modifie où et quand s'exécutent les tâches longues de la première
équipe, et c'est cela, et non les 5~\% de travail supplémentaire, qui fixe sa fin.

\section{Mise en cache, pannes et coût}

Lorsque chaque backtest recalcule ses variables (vingt minutes chacune), le balayage demande 6\,922 cœur-heures et 19.3~heures~; écrit sous forme de \omterm{def:pl:distributed-compute-and-schedulers:graph}{graphe de tâches}
avec cent étapes de variables, il demande 3\,514 cœur-heures et 11.4~heures~; le balayage suivant sur les mêmes données trouve les cent étapes dans le
cache et demande 3\,475 cœur-heures et 10.3~heures. Les pannes coûtent ce qui a été perdu~: une tentative échouée est reprise depuis son début, si bien qu'une panne
d'une tâche de sept heures coûte jusqu'à sept heures, ce qui explique que les tâches longues écrivent des points de reprise (livre~12, chapitre~23). Le coût d'un balayage est la capacité qu'il occupe~:
sur une grappe louée à l'heure-nœud, l'écart entre premier entré, premier sorti et le plus long d'abord avec copies de secours est de 235 heures-nœud pour
les mêmes résultats.

\section{Tutoriel~: un balayage sur une grappe simulée}\label{tut:pl:distributed-compute-and-schedulers:tut}

\begin{tutorial}
\textbf{Objectif.} Calibrer un balayage sur de vrais backtests, l'ordonnancer selon quatre politiques sur une grappe simulée, et mesurer ce que changent les copies de secours,
les parts équitables et la mise en cache. \textbf{État final~:} le \cref{tab:pl:distributed-compute-and-schedulers:policies} et
la \cref{fig:pl:distributed-compute-and-schedulers:profile}.
\begin{tutsteps}
\item \textbf{Calibrer}~: \texttt{bench\_jobgraph.py} exécute 24 petits backtests via \texttt{LocalExecutor} avec un seul travailleur~;
\texttt{measured\_sigma()}.
\item \textbf{Balayage}~: \texttt{sweep(sigma)}, \texttt{cluster()}, \texttt{lower\_bound}.
\item \textbf{Politiques}~: \texttt{policies(sigma)}~; \texttt{profile(sigma)} pour les courbes de cœurs occupés.
\item \textbf{Équipes}~: \texttt{fair\_share(sigma)}.
\item \textbf{Mise en cache}~: \texttt{caching(sigma)} sur le \omterm{def:pl:distributed-compute-and-schedulers:graph}{graphe de tâches}.
\end{tutsteps}
\textbf{Pour aller plus loin.} Donner à l'ordonnanceur des estimations parfaites et voir à quelle distance de la borne arrive le plus long d'abord~; ne lui fournir aucune estimation
(chaque tâche valant la médiane) et voir ce qui subsiste de son avantage.
\end{tutorial}

\section{Construction~: le graphe de tâches}\label{bld:pl:distributed-compute-and-schedulers:jobgraph}

\begin{build}
\bfield{Rôle} Une bibliothèque de \omterm{def:pl:distributed-compute-and-schedulers:graph}{graphes de tâches} et d'ordonnancement que la plateforme utilise pour planifier et exécuter les balayages, et un simulateur pour choisir les politiques
avant de dépenser des heures-nœud.
\bfield{Interface} \texttt{Task}, \texttt{Cluster}, \texttt{simulate(tasks, cluster, policy, spec, cache) -> Schedule}~;
\texttt{FIFO}, \texttt{LPT}, \texttt{FairShare}, \texttt{WorkStealing}~; \texttt{lower\_bound}~; \texttt{utilisation}~;
\texttt{LocalExecutor}.
\bfield{Règles} L'ordonnanceur ne voit que des estimations, jamais les durées~; une tentative échouée ou préemptée est relancée depuis son début~; une copie de secours s'exécute sur
un autre nœud, la copie perdante est interrompue et comptée comme gaspillage~; un succès de cache ne prend aucun temps~; tout est initialisé par une graine.
\bfield{Tests d'acceptation} \texttt{code/firm/jobgraph/tests/}~: l'exemple de Graham ($3,3,2,2,2$ sur deux cœurs~: 7 contre un optimum de 6)~; premier
entré, premier sorti face au plus long d'abord~; dépendances et borne de la chaîne~; partage équitable des cœurs~; \omterm{def:pl:distributed-compute-and-schedulers:stealing}{vol de tâches} sur une répartition initiale déséquilibrée~;
les copies de secours sauvant un nœud lent~; les relances~; les succès de cache~; l'ordre de l'exécuteur local.
\bfield{Pour aller plus loin} Classes de mémoire et tâches multicœurs avec backfilling~; capacité préemptible (spot) à prix réduit~; consommation qui s'estompe au fil du temps
dans l'\omterm{def:pl:distributed-compute-and-schedulers:fairshare}{ordonnancement à parts équitables}.
\end{build}

\begin{omsources}
\item R.~L. Graham, ``Bounds on multiprocessing timing anomalies'', \emph{SIAM Journal on Applied Mathematics} 17(2), 1969.
\item J.~Dean et S.~Ghemawat, ``MapReduce: simplified data processing on large clusters'', OSDI 2004.
\item J.~Dean et L.~A. Barroso, ``The tail at scale'', \emph{Communications of the ACM} 56(2), 2013.
\item R.~D. Blumofe and C.~E. Leiserson, ``Scheduling multithreaded computations by work stealing'', \emph{Journal of the ACM} 46(5), 1999.
\item Documentation de Slurm, ``Fair Tree Fairshare Algorithm''.
\end{omsources}

\section{Exercices}

\begin{exercise}[$\star$]\label{exo:pl:distributed-compute-and-schedulers:1}
Cinq tâches de durées 3, 3, 2, 2, 2 s'exécutent sur deux cœurs. Calculer la \omterm{def:pl:distributed-compute-and-schedulers:makespan}{durée totale} (makespan) avec le plus long d'abord, puis la meilleure durée possible.
\end{exercise}

\begin{exercise}[$\star$]\label{exo:pl:distributed-compute-and-schedulers:2}
La tâche médiane du balayage dure quinze minutes. Calculer la durée des quarante tâches longues sur un nœud normal, puis sur le nœud lent.
\end{exercise}

\begin{exercise}[$\star$]\label{exo:pl:distributed-compute-and-schedulers:3}
Expliquer pourquoi la borne inférieure du balayage du chapitre est le travail total divisé par le nombre de cœurs, et non la tâche la plus longue.
\end{exercise}

\begin{exercise}[$\star\star$]\label{exo:pl:distributed-compute-and-schedulers:4}
Expliquer pourquoi le seuil de copie de secours, égal à 2.5 fois l'estimation, se déclenche rarement pour une tâche saine, et ce que produirait un seuil de 1.5.
\end{exercise}

\begin{exercise}[$\star\star$]\label{exo:pl:distributed-compute-and-schedulers:5}
Expliquer pourquoi le \omterm{def:pl:distributed-compute-and-schedulers:stealing}{vol de tâches} se termine plus tard que le plus long d'abord, alors qu'aucun cœur n'est inactif tant qu'il reste du travail en file d'attente.
\end{exercise}

\begin{exercise}[$\star\star$]\label{exo:pl:distributed-compute-and-schedulers:6}
Le partage équitable permet à la seconde équipe de finir plus tôt~; expliquer ce qu'il change pour la fin de la première équipe, et pourquoi.
\end{exercise}

\begin{exercise}[$\star\star\star$]\label{exo:pl:distributed-compute-and-schedulers:7}
\emph{Programmation.} Lancer \texttt{policies} avec les tâches longues soumises en premier plutôt qu'en dernier, et expliquer ce qui arrive au premier entré, premier sorti et au plus long d'abord.
le plus long d'abord~?
\end{exercise}

\begin{exercise}[$\star\star\star$]\label{exo:pl:distributed-compute-and-schedulers:8}
\emph{Trouver la faille.} «~Le balayage a tourné toute la nuit sur une grappe de 384 cœurs et n'a consommé que 3\,649 heures-cœur, donc la grappe est trop grande~: nous
allons la réduire de moitié.~»
\end{exercise}

\section{Problème~: le balayage terminé le samedi}

\begin{problem}[{Problème du week-end --- dix mille backtests et un nœud lent}]\label{pb:pl:distributed-compute-and-schedulers:1}
Le balayage, la \omterm{def:pl:distributed-compute-and-schedulers:cluster}{grappe de calcul} et les politiques du chapitre.

\medskip\noindent\textbf{Partie I --- La charge de travail.}
\begin{enumerate}
\item Pourquoi un \omterm{def:pl:distributed-compute-and-schedulers:sweep}{balayage de paramètres} est-il \omterm{def:pl:distributed-compute-and-schedulers:sweep}{trivialement parallèle} (embarrassingly parallel)~?
\item Comment a-t-on obtenu la dispersion des durées des tâches~?
\item Quel est le travail total, et quelle est la borne inférieure sur 384 cœurs~?
\item À quel rang de soumission se trouvent les quarante tâches longues, et pourquoi~?
\item Qu'est-ce qui fait d'une tâche un \omterm{def:pl:distributed-compute-and-schedulers:makespan}{traînard} (straggler) ici, de trois façons~?
\end{enumerate}

\medskip\noindent\textbf{Partie II --- Politiques.}
\begin{enumerate}[resume]
\item Que garantit l'analyse de Graham pour l'ordonnancement en liste et pour le plus long d'abord~?
\item De quoi le plus long d'abord a-t-il besoin que le premier entré, premier sorti n'exige pas~?
\item Pourquoi le \omterm{def:pl:distributed-compute-and-schedulers:stealing}{vol de tâches} ne corrige-t-il pas les tâches longues~?
\item Que change le partage équitable entre deux équipes~?
\item Quand l'\omterm{def:pl:distributed-compute-and-schedulers:spec}{exécution spéculative} aide-t-elle, et quand gaspille-t-elle~?
\end{enumerate}

\medskip\noindent\textbf{Partie III --- Les chiffres.}
\begin{enumerate}[resume]
\item Donner la \omterm{def:pl:distributed-compute-and-schedulers:makespan}{durée totale} de chaque politique avec et sans copies de secours.
\item Pendant combien de temps la grappe est-elle restée presque inactive avec le premier entré, premier sorti~?
\item À quelle distance de la borne inférieure se trouve le meilleur ordonnancement, et pourquoi~?
\item Que fait le partage équitable pour la première moitié des tâches de la seconde équipe~?
\item Qu'est-ce que le \omterm{def:pl:distributed-compute-and-schedulers:graph}{graphe de tâches} et le cache permettent d'économiser~?
\end{enumerate}

\medskip\noindent\textbf{Partie IV --- Le verdict.}
\begin{enumerate}[resume]
\item Énoncer le \emph{résultat nommé}~: la \omterm{def:pl:distributed-compute-and-schedulers:makespan}{durée totale} et les heures-nœud du balayage sous chaque politique, le gain de l'\omterm{def:pl:distributed-compute-and-schedulers:spec}{exécution spéculative}, et
la distance du meilleur ordonnancement à la borne inférieure.
\item Quelles modifications apporter à l'ordonnancement par défaut des balayages de la plateforme~?
\item Quelles informations demander aux utilisateurs pour chaque balayage~?
\item Comment dimensionner la grappe~?
\item En une phrase~: qu'est-ce qui détermine la fin d'un balayage~?
\end{enumerate}
\end{problem}

\section{Questions d'entretien}

\begin{interviewq}[$\star$ \iqroles{développeur}]\label{iq:pl:distributed-compute-and-schedulers:1}
La dernière fraction de 1~\% des tâches de votre balayage prend autant de temps que les 99~\% premières. Que faites-vous~?
\end{interviewq}

\begin{interviewq}[$\star\star$ \iqroles{développeur}]\label{iq:pl:distributed-compute-and-schedulers:2}
Qu'est-ce que l'\omterm{def:pl:distributed-compute-and-schedulers:spec}{exécution spéculative}, et quand est-ce une mauvaise idée~?
\end{interviewq}

\begin{interviewq}[$\star\star$ \iqroles{développeur, chercheur}]\label{iq:pl:distributed-compute-and-schedulers:3}
Deux équipes partagent une grappe~; l'une soumet d'énormes balayages. Comment faites-vous pour que l'autre reste productive~?
\end{interviewq}

\begin{interviewq}[$\star\star$ \iqroles{chercheur}]\label{iq:pl:distributed-compute-and-schedulers:4}
Comment structureriez-vous dix mille backtests qui partagent des calculs de variables coûteux~?
\end{interviewq}

\begin{interviewq}[$\star\star\star$ \iqroles{développeur}]\label{iq:pl:distributed-compute-and-schedulers:5}
Concevez la couche d'ordonnancement de la \omterm{def:pl:distributed-compute-and-schedulers:cluster}{grappe de calcul} d'une plateforme de recherche.
\end{interviewq}
